% I&C 短縮実験：本論文から保留した補助的構造結果の元ソース。
% このファイルは短縮版から \input しない。数学的証明は保全。

\subsection{自明型付けを超える有限型付け}
\label{subsec:finite-typing-separations}

第~\ref{sec:prelim}節のフィルタに関する結果は，
十分条件による有用な構成法を与える．
しかし有限型付けは，
自明型付けクラスへ正則フィルタを後から適用するだけではない．
次の線形言語の例は，より強い分離を与えると同時に，
ファンアウトの鋭い効果を示す．

$\Sigma_{\mathrm{cd}}
:=\{\mathtt a,\mathtt b,\mathtt c,\mathtt d\}$
上の言語
\[
 L_{\mathrm{cd}}
 :=\{\mathtt a^n\mathtt c\mathtt b^n:n\ge0\}
 \cup
 \{\mathtt a^n\mathtt d\mathtt b^{2n}:n\ge0\}
\]
を考える．
$M_{\mathrm{cd}}
:=\{1,\eta_{\mathtt c},\eta_{\mathtt d},0\}$
とし，$1$ は単位元，$0$ は零元で，
$\{\eta_{\mathtt c},\eta_{\mathtt d},0\}$
から二つの元を取った積は常に $0$ とする．
準同型を
\[
 h_{\mathrm{cd}}(\mathtt a)
 =h_{\mathrm{cd}}(\mathtt b)=1,\qquad
 h_{\mathrm{cd}}(\mathtt c)=\eta_{\mathtt c},
 \qquad
 h_{\mathrm{cd}}(\mathtt d)=\eta_{\mathtt d}
\]
で定義する．

\begin{proposition}[ファンアウトに敏感な型付き言語]
\label{prop:strong-fanout-separation}
言語 $L_{\mathrm{cd}}$ は線形であり，
\(L_{\mathrm{cd}}\in\Cmcf{1}{h_{\mathrm{cd}}}\)
を満たす．
また任意の $k,\ell\ge0$ に対し
\(L_{\mathrm{cd}} \notin\Cmcf{1}{h_{\Sigma_{\mathrm{cd}}}^{k,\ell}}\)
である．
さらに
$L_{\mathrm{cd}}=L_0\cap L_{\mathrm{reg}}$
を満たす
$L_0\in\YSub{1}$ と正則言語 $L_{\mathrm{reg}}$
は存在しない．
しかも任意の有限モノイド準同型
$g:\Sigma_{\mathrm{cd}}^*\to N$ に対して
\(L_{\mathrm{cd}}\notin\Cmcf{2}{g}\)
となる．
\end{proposition}

\begin{proof}
線形文法
\[
 S\to C\mid D,\qquad
 C\to\mathtt c\mid\mathtt a C\mathtt b,\qquad
 D\to\mathtt d\mid\mathtt a D\mathtt b\mathtt b
\]
が $L_{\mathrm{cd}}$ を生成する．

$(1,h_{\mathrm{cd}})$-代入可能性について，
同じ型の非空因子 $x,y$ が
受理文脈を共有するとする．
型 $1$ の語には中央の記号が含まれないので，
文脈によって枝と左右の位置が固定され，
個数制約から $x=y$ となる．
型 $\eta_{\mathtt c}$ では
\(x=\mathtt a^{n_1}\mathtt c\mathtt b^{n_2}, \quad y=\mathtt a^{n_1'}\mathtt c\mathtt b^{n_2'}\)
と書ける．
共通文脈で受理されることから
$n_1-n_2=n_1'-n_2'$ が従う．
これは他のすべての $\mathtt c$-文脈でも不変な量である．
型 $\eta_{\mathtt d}$ では
$2n_1-n_2=2n_1'-n_2'$ が不変となる．
型 $0$ の分布は空である．
したがって
$L_{\mathrm{cd}}\in\Cmcf{1}{h_{\mathrm{cd}}}$ である．

$k,\ell$ を固定し，
$m\ge\max\{1,k+\ell\}$ を選ぶ．
\(x=\mathtt a^m\mathtt c\mathtt b^m,\qquad y=\mathtt a^m\mathtt d\mathtt b^{2m}\)
は同じ $h_{\Sigma_{\mathrm{cd}}}^{k,\ell}$-型をもち，
空文脈を共有する．
しかし
$\mathtt a\blank_1\mathtt b$ は
$x$ を受理する一方，$y$ を拒否する．
よって，すべての接頭辞--接尾辞型付けは
ファンアウト1でも失敗する．

$L_{\mathrm{cd}}=L_0\cap L_{\mathrm{reg}}$ で，
$L_0\in\YSub{1}$，
$L_{\mathrm{reg}}$ が正則だと仮定する．
$\mathtt c,\mathtt d\in L_0$ なので，
代入可能性から
$\D_{L_0}(\mathtt c)=\D_{L_0}(\mathtt d)$ が従う．
さらに
$\mathtt a\mathtt d\mathtt b^2,
 \mathtt a\mathtt c\mathtt b\in L_0$
を使うと
$\D_{L_0}(\mathtt c\mathtt b)
=\D_{L_0}(\mathtt c\mathtt b^2)$
が得られる．
したがって
\(\mathtt a^m\mathtt c\mathtt b^{m+t}\in L_0 \qquad(m\ge1,\ t\ge0)\)
である．
$L_{\mathrm{reg}}$ の DFA の状態数より $m$ を大きくし，
受理される標的語の末尾の
$\mathtt b^m$ ブロックをポンピングすると，
ある $t>0$ について同じ語が
$L_{\mathrm{reg}}$ にも入る．
これは仮定した交差表示に矛盾する．

最後に有限準同型
$g:\Sigma_{\mathrm{cd}}^*\to N$ に対し，
$1\le i<j$ を
\(g(\mathtt a^i)=g(\mathtt a^j),\qquad g(\mathtt b^i)=g(\mathtt b^j)\)
となるように選ぶ．
タプル
$(\mathtt a^i,\mathtt b^i)$ と
$(\mathtt a^j,\mathtt b^j)$
は成分別 $g$-型が等しく，
$\blank_1\mathtt c\blank_2$
によって両方とも受理される．
一方，
$\blank_1\mathtt d\blank_2\mathtt b^i$
は最初のタプルだけを受理する．
$i<j$ なら
$\mathtt a^j\mathtt d\mathtt b^{i+j}
\notin L_{\mathrm{cd}}$
だからである．
よってアリティ2の条件は，
任意の有限 $g$ に対して破れる．
\end{proof}

したがって $L_{\mathrm{cd}}$ は，
ファンアウト1におけるすべての接頭辞--接尾辞型付けと
任意の正則フィルタによる表現から有限型付けを分離する．
しかしアリティ2を検査するようにすると，
どの有限型付けでもこの言語を救えない．

\subsection{局所的な境界制御と任意の正則フィルタ}
\label{subsec:comparison-separations}

局所検査可能性の標準用語に従う
~\cite{McNaughtonPapert1971}．
$L_{\mathrm{sl}}\subseteq\Sigma^*$ が
\emph{厳密2局所（strictly 2-local）}であるとは，
ある集合
$\mathsf{Adj}
\subseteq(\Sigma\cup\{\vdash,\dashv\})^2$
が存在して，
\[
 w\in L_{\mathrm{sl}}
 \quad\Longleftrightarrow\quad
 \text{語 $\vdash w\dashv$ の隣接対がすべて
       $\mathsf{Adj}$ に属する}
\]
となることをいう．
ここで $\vdash,\dashv\notin\Sigma$ は固定した端点標識である．

\begin{lemma}[厳密2局所条件は境界型付きである]
\label{lem:sl2-boundary-control}
$L_{\mathrm{sl}}\subseteq\Sigma^*$ が厳密2局所なら，
任意の $f\ge1$ について
\(L_{\mathrm{sl}} \in\Cmcf{f}{h_\Sigma^{1,1}}\)
が成り立つ．
\end{lemma}

\begin{proof}
同じ $h_\Sigma^{1,1}$-型をもつ非空タプル
$\tuple x,\tuple y$ が，
受理する文脈 $E_0$ を共有するとする．
各 $y_i$ の内部の隣接対は
$E_0[\tuple y]$ に現れるので許容される．
$E[\tuple x]\in L_{\mathrm{sl}}$ のとき，
$E[\tuple y]$ の各隣接対は，
固定終端部分の内部，
ある $y_i$ の内部，
または成分の境界を横切る対のいずれかである．
それらは順に
$E[\tuple x]$，
$E_0[\tuple y]$，
あるいは境界型で符号化された先頭・末尾文字の一致
によって許容される．
端点標識を含む対は変化しない．
したがって $E[\tuple y]\in L_{\mathrm{sl}}$ であり，
対称性から主張が従う．
\end{proof}

\paragraph{標的依存の任意の正則フィルタ．}
比較のため，固定アルファベット上に
\[
 \mathcal F_f^{\mathrm{arb}}
 :=\bigl\{L_0\cap R:
       L_0\in\YSub{f},\ R\text{ は正則}\bigr\}
       \cap f\text{-}\mathrm{MCFL}
\]
と定義する．
このクラスはすべての正則言語を含む
（$\Sigma^*\in\YSub{f}$ とすればよい）．
よって Gold の障害により，
単一の正例学習クラスとしては
TxtEx 同定できない．
ここでは追加の学習可能クラスとしてではなく，
表現力を比較するための基準としてだけ用いる．

内在的ランク2の証拠言語はこの大きな断片に属する：
\(L_{\mathrm{intr}} =\Phi_{\mathrm{intr}}^{-1}(0)\cap Q_{\mathrm{intr}}.\)
したがって，この言語が示すのは内在的規則ランクであって，
任意の正則フィルタからの分離ではない．
命題~\ref{prop:strong-fanout-separation} は
ファンアウト1でその分離を示した．
次の定理は，型付き非文脈自由の
ファンアウト2・ランク1の標的について分離を示す．

\begin{theorem}[任意の正則フィルタを超えるファンアウト2の型付き言語]
\label{thm:arbitrary-filter-separation}
$\Sigma_\star:=\{\mathtt a,\mathtt b,\mathtt c,\mathtt d\}$ とし，
\[
 L_\star
 :=\{\mathtt a^n\mathtt c^{n+1}\mathtt b^n:n\ge0\}
   \cup
   \{\mathtt a^n\mathtt d^{n+1}\mathtt b^{2n}:n\ge0\}
\]
とおく．
準同型
$g_\star:\Sigma_\star^*
\to(\mathcal P(\{\mathtt c,\mathtt d\}),\cup)$ を
\[
 g_\star(\mathtt a)=g_\star(\mathtt b)=\varnothing,
 \qquad
 g_\star(\mathtt c)=\{\mathtt c\},
 \qquad
 g_\star(\mathtt d)=\{\mathtt d\}
\]
で定め，
\(h_\star:=(h_{\Sigma_\star}^{1,1},g_\star)\)
とおく．
このとき $L_\star$ は非文脈自由で，
最小ファンアウトは2，
ファンアウト2での最小規則ランクは1である．
さらに任意の $f\ge2$ について
\(L_\star \in\Cmcf{f}{h_\star} \setminus\mathcal F_f^{\mathrm{arb}}\)
が成り立つ．
同値に，
$L_0\in\YSub{f}$ と正則言語 $R_{\mathrm{reg}}$ により
$L_\star=L_0\cap R_{\mathrm{reg}}$
と表すことはできない．
\end{theorem}

\begin{proof}
良形でファンアウト2・ランク1の文法を
\[
\begin{aligned}
 A_{\mathtt c}
   &\to(\mathtt a\mathtt c,\mathtt b)
     \mid(\mathtt a\xi_1^1\mathtt c,\xi_1^2\mathtt b)(A_{\mathtt c}),\\
 A_{\mathtt d}
   &\to(\mathtt a\mathtt d,\mathtt b\mathtt b)
     \mid(\mathtt a\xi_1^1\mathtt d,\xi_1^2\mathtt b\mathtt b)(A_{\mathtt d}),\\
 S&\to\mathtt c\mid\mathtt d
   \mid(\xi_1^1\mathtt c\xi_1^2)(A_{\mathtt c})
   \mid(\xi_1^1\mathtt d\xi_1^2)(A_{\mathtt d})
\end{aligned}
\]
で与える．
これは正確に $L_\star$ を生成する．
正則言語
$\mathtt a^*\mathtt c^+\mathtt b^*$
との交差は
$\{\mathtt a^n\mathtt c^{n+1}\mathtt b^n:n\ge0\}$
であり，これは文脈自由でない．
閉包性を使う短い証明のため，
新しい記号 $\mathtt\#$ を導入し，
$\mathtt a,\mathtt b,\mathtt c$ は固定して
$\mathtt\#$ を $\mathtt c$ に写す準同型を考える．
この枝の逆像を
$\mathtt a^*\mathtt c^*\mathtt\#\mathtt b^*$
と交差すると，
\(\{\mathtt a^n\mathtt c^n\mathtt\#\mathtt b^n:n\ge0\}\)
となる．
$\mathtt\#$ を消去すれば
$\{\mathtt a^n\mathtt c^n\mathtt b^n:n\ge0\}$
が文脈自由となってしまい，矛盾する．
よって最小ファンアウトは2である．
言語が無限であることと上に示した文法から，
ファンアウト2における最小規則ランクは1である．

型付き代入可能性を示すため，
\[
 Q_\star:=
   \mathtt a^*\mathtt c^+\mathtt b^*
   \cup\mathtt a^*\mathtt d^+\mathtt b^*,
 \qquad
 \Psi(w):=(|w|_{\mathtt a},|w|_{\mathtt b},
            |w|_{\mathtt c},|w|_{\mathtt d})
\]
とおく．
厳密2局所言語 $Q_\star$ は
補題~\ref{lem:sl2-boundary-control} により境界型付きである．
同じ $h_\star$-型をもつタプル
$\tuple x,\tuple y$ が
受理文脈 $E_0$ を共有するとし，
\(\Delta:=\sum_i\Psi(y_i)-\sum_i\Psi(x_i)\)
とおく．
両者の $g_\star$-型が等しいので，
$E_0[\tuple x]$ と $E_0[\tuple y]$
は同じ枝に属する．
したがってある $q\in\mathbb Z$ に対し，
その枝が $\mathtt c$ なら
$\Delta=q(1,1,1,0)$，
$\mathtt d$ なら
$\Delta=q(1,2,0,1)$ である．

$F[\tuple x]\in L_\star$ とする．
境界型付けにより
$F[\tuple y]\in Q_\star$ であり，
$g_\star$-型の一致から
二つの充填結果は同じ $F$-枝に属する．
$E_0$ と $F$ が同じ枝を選ぶなら，
上の個数差はその枝のアフィンな関係を保存する．
別の枝を選ぶ場合，
いずれのタプル成分にも
$\mathtt c,\mathtt d$ は含まれない．
そのため $\Delta$ の中央の二座標は0となり，
$q=0$ が強制される．
いずれの場合も $F[\tuple y]\in L_\star$ であり，
対称性から
$L_\star\in\Cmcf{f}{h_\star}$ が従う．

最後に
$L_\star=L_0\cap R_{\mathrm{reg}}$，
$L_0\in\YSub{f}$，
$R_{\mathrm{reg}}$ は正則，と仮定する．
$\mathtt c,\mathtt d\in L_0$ なので，
アリティ1の代入可能性から
$\D_{L_0}(\mathtt c)=\D_{L_0}(\mathtt d)$
となる．
よって
$\mathtt a\mathtt d^2\mathtt b^2\in L_0$
で二つの $\mathtt d$ を順に置き換えると，
まず
$\mathtt a\mathtt c\mathtt d\mathtt b^2\in L_0$，
次に
$\mathtt a\mathtt c^2\mathtt b^2\in L_0$
が得られる．
これと
$\mathtt a\mathtt c^2\mathtt b\in L_0$
から
\(\D_{L_0}(\mathtt c^2\mathtt b) =\D_{L_0}(\mathtt c^2\mathtt b^2)\)
が従う．
したがって任意の $m\ge1,t\ge0$ について
$\mathtt a^m\mathtt c^{m+1}\mathtt b^{m+t}\in L_0$
となる．
$m$ を $R_{\mathrm{reg}}$ の DFA の状態数以上にすると，
受理される標的語の末尾の
$\mathtt b^m$ ブロックをポンピングでき，
ある $t>0$ について
同じ語が $R_{\mathrm{reg}}$ に属する．
この語は交差には属するが
$L_\star$ には属さないため，矛盾する．
\end{proof}

よって単一の標的 $L_\star$ は，
Yoshinaka の言語にどのような有限状態フィルタを
後から適用しても表現できないことを示す．
この結論は，フィルタに用いる型付けが
$h_\star$ をどれほど細かく精密化しても変わらない．

このランク1の分離は，
定理~\ref{thm:intrinsic-ranktwo-witness} を補完する．
後者は，すべての有限ファンアウトにわたり
最小規則ランク2をもつ
境界型付き非文脈自由標的を与える．

\begin{proposition}[固定窓では括弧付きCFG族のファンアウト2代入可能性を回復できない]
\label{prop:fixed-window-doubling-fanout-two}
$\Gamma_{\mathrm{br}}=\{\mathtt a,\lfloor,\rfloor,\lceil,\rceil\}$ とし，
関連CFG原稿~\cite{kuriyamaCFG} の言語族 $L_n$ を
\[
 E_0\to\mathtt a,\qquad
 E_i\to\lfloor E_{i-1}E_{i-1}\rfloor
       \mid\lceil E_{i-1}\rceil\quad(1\le i\le n)
\]
の開始非終端記号 $E_n$ から生成する．
$h_{\mathrm U}$ は $\lceil,\rceil$ を2元半束の吸収元に，
それ以外の文字を単位元に写す準同型とする．
任意の $k,\ell\ge0$ および
$n\ge\max\{3,k+\ell+1\}$ について
\[
 L_n\notin\Cmcf{2}{h_{\Gamma_{\mathrm{br}}}^{k,\ell}},
 \qquad
 L_n\notin\Cmcf{2}{(h_{\Gamma_{\mathrm{br}}}^{k,\ell},h_{\mathrm U})}.
\]
したがって，$h_{\mathrm U}$ と組み合わせても，
\emph{固定}接頭辞--接尾辞窓を加えるだけでは
この族のファンアウト2での代入可能性は回復しない．
\end{proposition}
\begin{proof}
$W_0=\mathtt a$，
$W_i=\lfloor W_{i-1}W_{i-1}\rfloor$ とおく．
$L_n$ のうち天井括弧を含まない唯一の語は $W_n$ である．
完全二分木の語 $W_{n-1}$ を
$W_{n-1}=B\mathtt a R$，
$B=\lfloor^{n-1}$ と書けば
\[
 W_n=\lfloor B\mathtt a R\,B\mathtt a R\rfloor
\]
となる．$B,R$ はいずれも非空で，
$|B|=n-1\ge k+\ell$，
$|R|=3\cdot2^{n-1}-n-2\ge n-1\ge k+\ell$ である．

次の2タプルと2つの文文脈を考える：
\[
 \begin{aligned}
 \tuple x&=(B,\ R B\mathtt a R),&
 \tuple y&=(B\mathtt a R B,\ R),\\
 E&=\lfloor\blank_1\mathtt a\blank_2\rfloor,&
 F&=\blank_1\lfloor\mathtt a\blank_2\rfloor.
 \end{aligned}
\]
どちらの文脈でも2つの穴の間の文字列は非空であり，
定義~\ref{def:tuple-substitutable} に適合する．
第1成分の $B$ と $B\mathtt a R B$ は，
長さ $k$ の接頭辞と長さ $\ell$ の接尾辞が一致する．
第2成分の $R B\mathtt a R$ と $R$ も同様である．
各成分は $k+\ell$ 以上の長さを持ち，
天井括弧を含まない．
よって $\tuple x,\tuple y$ の成分別の型は，
主張にある両方の型付けで一致する．
しかし
\[
 E[\tuple x]=E[\tuple y]=F[\tuple x]=W_n,
 \qquad
 F[\tuple y]=B\mathtt a R B\lfloor\mathtt a R\rfloor
 \ne W_n.
\]
$F[\tuple y]$ の最初の $\mathtt a$ の前には
$\lfloor$ が $n-1$ 個しかないのに対し，
$W_n$ では $n$ 個あるためである．
$F[\tuple y]$ にも天井括弧は含まれないため，
$F[\tuple y]\notin L_n$ である．
従って両方の型付けについて
$(2,h)$-タプル代入可能性が破れる．
\end{proof}

\begin{remark}
本命題は元のCFG族が高ファンアウトの標的クラスに
\emph{属さない}ことを示すのであって，
ファンアウト2における特徴データ長の超多項式下界ではない．
特に問題~\ref{op:fanout-two} は未解決であり，
その解決には $\Cmcf{2}{h}$ に実際に属する
\emph{別の}言語族の構成が必要である．
\end{remark}



\subsection{希少な型と2成分の局所的可動性}
\label{subsec:rare-type-mobility}

\begin{lemma}[$h$ の型全体の削除]
\label{lem:mcfg-delete-type}
有限半群型付け $s:\Sigma^+\twoheadrightarrow\mathsf S$，
$f\ge1$，$L\in\Cmcf{f}{s}$，$m\in\mathsf S$ に対し，
$L^{\setminus m}:=L\setminus s^{-1}(m)$ も
$\Cmcf{f}{s}$ に属する．
\end{lemma}
\begin{proof}
$h=\widehat s$，$M=\mathsf S^\sharp$ とおく．
$h^{-1}(M\setminus\{m\})$ は正規言語であり，
$(f,s)$-タプル代入可能でもある．実際，成分別の型が
一致すれば，どの文文脈でも全体の型が一致する．
共通部分は代入可能性を保存し
（命題~\ref{prop:abelian-coset-filter}），
$f$-MCFL は正規言語との共通部分で閉じる
~\cite{SekiEtAl1991}．
\end{proof}

\begin{proposition}[希少な型は長い正例を強制する]
\label{prop:mcfg-rare-type-sample}
有限半群型付け $s:\Sigma^+\twoheadrightarrow\mathsf S$ と
$L\in\Cmcf{f}{s}$，$m\in s(L\cap\Sigma^+)$ をとり，
$h=\widehat s$，
$\mu_m(L)=\min\{|w|:w\in L\cap\Sigma^+,\ s(w)=m\}$ とおく．
$\mathcal B^{\le r}_{f,h}$（$r\ge1$）または
$\mathcal B_{f,h}$ の任意の特徴標本には
型 $m$ の語が含まれ，従って符号化総長は
$\mu_m(L)+1$ 以上である．
より一般に，$L$ と $L^{\setminus m}$ の両方に
特徴標本を持つ任意の集合駆動型再構成器についても
同じ必要性が成立する．
\end{proposition}
\begin{proof}
型 $m$ を避ける標本 $K\subseteq L$ は
$K\subseteq L^{\setminus m}$ を満たす．
補題~\ref{lem:mcfg-delete-type} と
命題~\ref{prop:reconstruction-soundness} を
標的 $L^{\setminus m}$ に適用すると，
$K$ からの再構成言語は真部分言語
$L^{\setminus m}$ に含まれる．
一般の集合駆動型再構成器については，
$L$ の特徴標本 $\mathcal F$ が型 $m$ を避けると仮定し，
$L^{\setminus m}$ の特徴標本を $\mathcal F'$ とする．
有限集合 $\mathcal F\cup\mathcal F'$ は両標的に含まれ，
双方の特徴標本を含むので，同じ再構成結果が
異なる二言語に等しくなり矛盾する．
\end{proof}

\begin{lemma}[ファンアウト2の冪移送]
\label{lem:mcfg-power-transfer}
$f\ge2$、$L\in\Cmcf{f}{h}$、$u\in\Sigma^+$ とする．
$g=h(u)$ とし，ある $i_0,\pi\ge1$ が
全ての $i\ge i_0$ について
$g^{i+\pi}=g^i$ を満たすとする．
$s,r\ge i_0$、$d\in\pi\mathbb Z_{>0}$ とする．
$L$ のある語が $u^{s+d+r+2}$ を因子に持つなら，
任意の $X,Z\in\Sigma^*$、$Y\in\Sigma^+$ に対し
\[
 Xu^sYu^{d+r}Z\in L
 \quad\Longleftrightarrow\quad
 Xu^{s+d}Yu^rZ\in L
\]
が成り立つ．
\end{lemma}
\begin{proof}
受理語を $Pu^{s+d+r+2}Q$ と書く．
$\tuple x=(u^s,u^{d+r})$、
$\tuple y=(u^{s+d},u^r)$ とすると，
両者の成分別 $h$-型は一致する．
文脈 $E=Pu\blank_1u\blank_2Q$ は内部区切りが非空で，
$E[\tuple x]=E[\tuple y]=Pu^{s+d+r+2}Q\in L$ である．
よって2成分代入可能性から
両タプルの分布は一致する．
任意の適切な文脈 $X\blank_1Y\blank_2Z$
に適用すれば主張が従う．
\end{proof}

\begin{lemma}[有界な局所シフト対]
\label{lem:mcfg-local-shift}
有限半群型付け $s:\Sigma^+\to\mathsf S$ をとり，
$\Sigma$ は空でないとする．$h=\widehat s$，
$M=\mathsf S^\sharp$，$c=|M|^2|\Sigma|$ とし，$R_c(5)$ を
5頂点完全グラフに対する $c$ 色Ramsey数とする．
$|w|\ge R_c(5)$ を満たす任意の語は
$w=P\,x\,p\,q\,p\,y\,Q$ と分解でき，
$x,y\in\Sigma^+$、$p\in\Sigma$、
$q\in\Sigma^*$ および
\[
 h(x)=h(xpq),\qquad h(qpy)=h(y)
\]
を満たす．従って $w\in L\in\Cmcf{2}{h}$ ならば
$(x,qpy)$ と $(xpq,y)$ の分布は一致する．
\end{lemma}
\begin{proof}
最初の $R_c(5)$ 個のカット位置から $i<j$ を選び，
$(h(w[i:j]),h(w[i+1:j]),w[i])$ で辺を彩色する．
単色5点集合 $i_1<\cdots<i_5$ を取り，
$x=w[i_1:i_2]$、
$p=w[i_2]=w[i_3]$、
$q=w[i_2+1:i_3]$、
$y=w[i_3+1:i_5]$ とする．
$P,Q$ は残りの接頭辞・接尾辞とする．
辺 $(i_1,i_2),(i_1,i_3)$ の第一座標が
$h(x)=h(xpq)$ を与え，
$(i_2,i_5),(i_3,i_5)$ の第二座標が
$h(qpy)=h(y)$ を与える．
内部区切り $p$ が非空である文脈
$P\blank_1p\blank_2Q$ は両タプルを受理するため，
2成分代入可能性によって分布が一致する．
\end{proof}

\begin{theorem}[任意の固定有限型付けは括弧付き倍化族を救済できない]
\label{thm:mcfg-all-typing-doubling-obstruction}
任意の有限半群準同型
$s:\Gamma_{\mathrm{br}}^+\to\mathsf S$ について，
命題~\ref{prop:fixed-window-doubling-fanout-two} の $L_n$ は
\[
 n\ge2|\mathsf S|+3\quad\Longrightarrow\quad
 L_n\notin\Cmcf{2}{s}
\]
を満たす．$f\ge2$ に対しても
$L_n\notin\Cmcf{f}{s}$ となる．
\end{theorem}
\begin{proof}
$g=s(\lfloor)$ の冪の指数を $i_0\ge1$、
周期を $\pi\ge1$ とし，
$i_0+\pi\le|\mathsf S|+1$、$i_0\le|\mathsf S|$ とする．
$n\ge2i_0+\pi+2$ の場合を考え，
$a=i_0$、$d=\pi$、
$r=n-a-d-2\ge i_0$ とおく．
$W_n=\lfloor W_{n-1}W_{n-1}\rfloor$ は
$L_n$ に属する唯一の天井括弧なしの語で，
$\lfloor^n$ から始まる．
ゆえに $u=\lfloor$ として
補題~\ref{lem:mcfg-power-transfer} の前提が成立する．

同じ $W_n$ の先頭ランから，
最初の $\lfloor$ を除いた位置に
第一ブロック $\lfloor^a$ を取る．
第二ブロック $\lfloor^{d+r}$ は
二番目の $W_{n-1}$ の先頭ラン内に取る．
両者の間の文字列 $Y$ は最初の $\mathtt a$ を
含むため非空である．従って適切な $X,Z$ により
$W_n=X\lfloor^aY\lfloor^{d+r}Z$ と書ける．
冪移送から
$X\lfloor^{a+d}Y\lfloor^rZ\in L_n$
が従うはずだが，これは天井括弧を含まず，
最初の $\mathtt a$ の前の $\lfloor$ が
$n+d$ 個ある．一方，唯一の天井括弧なしの語
$W_n$ では $n$ 個であり，矛盾する．
$2i_0+\pi+2\le2|\mathsf S|+3$ より一様な閾値が得られる．
\end{proof}

\begin{remark}
定理~\ref{thm:mcfg-all-typing-doubling-obstruction} は，
命題~\ref{prop:fixed-window-doubling-fanout-two} の
窓型付けに限った障害を任意の有限型付けへ強化する．
窓型付け固有の鋭い閾値は旧命題に残す．
これは言語の\emph{クラス非所属}の定理であり，
問題~\ref{op:fanout-two} が問う特徴データの
高ファンアウト下界ではない．
\end{remark}

窓型付けと任意の有限群を組み合わせた場合の拡張は，
\ref{app:window-group} で証明した．
可換群の場合のより小さい定数は
\ref{app:window-abelian} に残す．
固定窓のみの場合の二次上界と窓・有限群型付けの場合の
三次上界について，文法サイズの指数が最適かどうかは未解決である．
非可換群の定数 $M_{A,f}$ の $f$ への二重指数的依存を下げられるかも今後の課題である．
付録~\ref{app:local-group} は，正像が局所群となるすべての型付け，
すなわち2元半束を含まない最大のクラスに，この元文法厚みによる制御を拡張する．
対応する下界はファンアウト1でのみ知られている（系~\ref{cor:pseudovariety-dichotomy}）．
系~\ref{cor:mcfg-li-group-polydata} は，
$\mathbf{LI}\vee\mathbf{Grp}$ による見通しのよい十分条件を与える．
より広い局所群型付けと擬多様体レベルの二分法は
付録~\ref{app:local-group} で別に扱う．
高ファンアウトの下界問題は引き続き未解決である．

今後の方向性としては，
複数のファンアウトバイアスにまたがる適応的学習，
値域サイズの事前上界なしでの有用な型付けの学習，
型付きクラス内部の規則ランク1部分クラスの特徴づけ，
全多重文脈に基づく合同的 MCFG との関係の解明，
および PMCFG 型のコピーを許す再構成への拡張がある．

